\documentclass[a4paper,12pt,twoside]{memoir}

% Castellano
\usepackage[spanish,es-tabla]{babel}
\selectlanguage{spanish}
\usepackage[utf8]{inputenc}
\usepackage[T1]{fontenc}
\usepackage{lmodern} % Scalable font
\usepackage{microtype}
\usepackage{placeins}

\RequirePackage{booktabs}
\RequirePackage[table]{xcolor}
\RequirePackage{xtab}
\RequirePackage{multirow}

% Links
\PassOptionsToPackage{hyphens}{url}\usepackage[colorlinks]{hyperref}
\hypersetup{
	allcolors = {blue}
}

% Ecuaciones
\usepackage{amsmath}

% Rutas de fichero / paquete
\newcommand{\ruta}[1]{{\sffamily #1}}

% Párrafos
\nonzeroparskip

% Huérfanas y viudas
\widowpenalty100000
\clubpenalty100000

% Imágenes

% Comando para insertar una imagen en un lugar concreto.
% Los parámetros son:
% 1 --> Ruta absoluta/relativa de la figura
% 2 --> Texto a pie de figura
% 3 --> Tamaño en tanto por uno relativo al ancho de página
\usepackage{graphicx}
\newcommand{\imagen}[3]{
	\begin{figure}[!h]
		\centering
		\includegraphics[width=#3\textwidth]{#1}
		\caption{#2}\label{fig:#1}
	\end{figure}
	\FloatBarrier
}

% Comando para insertar una imagen sin posición.
% Los parámetros son:
% 1 --> Ruta absoluta/relativa de la figura
% 2 --> Texto a pie de figura
% 3 --> Tamaño en tanto por uno relativo al ancho de página
\newcommand{\imagenflotante}[3]{
	\begin{figure}
		\centering
		\includegraphics[width=#3\textwidth]{#1}
		\caption{#2}\label{fig:#1}
	\end{figure}
}

% El comando \figura nos permite insertar figuras comodamente, y utilizando
% siempre el mismo formato. Los parametros son:
% 1 --> Porcentaje del ancho de página que ocupará la figura (de 0 a 1)
% 2 --> Fichero de la imagen
% 3 --> Texto a pie de imagen
% 4 --> Etiqueta (label) para referencias
% 5 --> Opciones que queramos pasarle al \includegraphics
% 6 --> Opciones de posicionamiento a pasarle a \begin{figure}
\newcommand{\figuraConPosicion}[6]{%
  \setlength{\anchoFloat}{#1\textwidth}%
  \addtolength{\anchoFloat}{-4\fboxsep}%
  \setlength{\anchoFigura}{\anchoFloat}%
  \begin{figure}[#6]
    \begin{center}%
      \Ovalbox{%
        \begin{minipage}{\anchoFloat}%
          \begin{center}%
            \includegraphics[width=\anchoFigura,#5]{#2}%
            \caption{#3}%
            \label{#4}%
          \end{center}%
        \end{minipage}
      }%
    \end{center}%
  \end{figure}%
}

%
% Comando para incluir imágenes en formato apaisado (sin marco).
\newcommand{\figuraApaisadaSinMarco}[5]{%
  \begin{figure}%
    \begin{center}%
    \includegraphics[angle=90,height=#1\textheight,#5]{#2}%
    \caption{#3}%
    \label{#4}%
    \end{center}%
  \end{figure}%
}
% Para las tablas
\newcommand{\otoprule}{\midrule [\heavyrulewidth]}
%
% Nuevo comando para tablas pequeñas (menos de una página).
\newcommand{\tablaSmall}[5]{%
 \begin{table}
  \begin{center}
   \rowcolors {2}{gray!35}{}
   \begin{tabular}{#2}
    \toprule
    #4
    \otoprule
    #5
    \bottomrule
   \end{tabular}
   \caption{#1}
   \label{tabla:#3}
  \end{center}
 \end{table}
}

%
% Nuevo comando para tablas pequeñas (menos de una página).
\newcommand{\tablaSmallSinColores}[5]{%
 \begin{table}[H]
  \begin{center}
   \begin{tabular}{#2}
    \toprule
    #4
    \otoprule
    #5
    \bottomrule
   \end{tabular}
   \caption{#1}
   \label{tabla:#3}
  \end{center}
 \end{table}
}

\newcommand{\tablaApaisadaSmall}[5]{%
\begin{landscape}
  \begin{table}
   \begin{center}
    \rowcolors {2}{gray!35}{}
    \begin{tabular}{#2}
     \toprule
     #4
     \otoprule
     #5
     \bottomrule
    \end{tabular}
    \caption{#1}
    \label{tabla:#3}
   \end{center}
  \end{table}
\end{landscape}
}

%
% Nuevo comando para tablas grandes con cabecera y filas alternas coloreadas en gris.
\newcommand{\tabla}[6]{%
  \begin{center}
    \tablefirsthead{
      \toprule
      #5
      \otoprule
    }
    \tablehead{
      \multicolumn{#3}{l}{\small\sl continúa desde la página anterior}\\
      \toprule
      #5
      \otoprule
    }
    \tabletail{
      \hline
      \multicolumn{#3}{r}{\small\sl continúa en la página siguiente}\\
    }
    \tablelasttail{
      \hline
    }
    \bottomcaption{#1}
    \rowcolors {2}{gray!35}{}
    \begin{xtabular}{#2}
      #6
      \bottomrule
    \end{xtabular}
    \label{tabla:#4}
  \end{center}
}

%
% Nuevo comando para tablas grandes con cabecera.
\newcommand{\tablaSinColores}[6]{%
  \begin{center}
    \tablefirsthead{
      \toprule
      #5
      \otoprule
    }
    \tablehead{
      \multicolumn{#3}{l}{\small\sl continúa desde la página anterior}\\
      \toprule
      #5
      \otoprule
    }
    \tabletail{
      \hline
      \multicolumn{#3}{r}{\small\sl continúa en la página siguiente}\\
    }
    \tablelasttail{
      \hline
    }
    \bottomcaption{#1}
    \begin{xtabular}{#2}
      #6
      \bottomrule
    \end{xtabular}
    \label{tabla:#4}
  \end{center}
}

%
% Nuevo comando para tablas grandes sin cabecera.
\newcommand{\tablaSinCabecera}[5]{%
  \begin{center}
    \tablefirsthead{
      \toprule
    }
    \tablehead{
      \multicolumn{#3}{l}{\small\sl continúa desde la página anterior}\\
      \hline
    }
    \tabletail{
      \hline
      \multicolumn{#3}{r}{\small\sl continúa en la página siguiente}\\
    }
    \tablelasttail{
      \hline
    }
    \bottomcaption{#1}
  \begin{xtabular}{#2}
    #5
   \bottomrule
  \end{xtabular}
  \label{tabla:#4}
  \end{center}
}



\definecolor{cgoLight}{HTML}{EEEEEE}
\definecolor{cgoExtralight}{HTML}{FFFFFF}

%
% Nuevo comando para tablas grandes sin cabecera.
\newcommand{\tablaSinCabeceraConBandas}[5]{%
  \begin{center}
    \tablefirsthead{
      \toprule
    }
    \tablehead{
      \multicolumn{#3}{l}{\small\sl continúa desde la página anterior}\\
      \hline
    }
    \tabletail{
      \hline
      \multicolumn{#3}{r}{\small\sl continúa en la página siguiente}\\
    }
    \tablelasttail{
      \hline
    }
    \bottomcaption{#1}
    \rowcolors[]{1}{cgoExtralight}{cgoLight}

  \begin{xtabular}{#2}
    #5
   \bottomrule
  \end{xtabular}
  \label{tabla:#4}
  \end{center}
}



\graphicspath{ {./img/} }

% Capítulos
\chapterstyle{bianchi}
\newcommand{\capitulo}[2]{
	\setcounter{chapter}{#1}
	\setcounter{section}{0}
	\setcounter{figure}{0}
	\setcounter{table}{0}
	\chapter*{\thechapter.\enskip #2}
	\addcontentsline{toc}{chapter}{\thechapter.\enskip #2}
	\markboth{#2}{#2}
}

% Apéndices
\renewcommand{\appendixname}{Apéndice}
\renewcommand*\cftappendixname{\appendixname}

\newcommand{\apendice}[1]{
	%\renewcommand{\thechapter}{A}
	\chapter{#1}
}

\renewcommand*\cftappendixname{\appendixname\ }

% Formato de portada
\makeatletter
\usepackage{xcolor}
\newcommand{\tutor}[1]{\def\@tutor{#1}}
\newcommand{\course}[1]{\def\@course{#1}}
\definecolor{cpardoBox}{HTML}{E6E6FF}
\def\maketitle{
  \null
  \thispagestyle{empty}
  % Cabecera ----------------
\noindent\includegraphics[width=\textwidth]{cabecera}\vspace{1cm}%
  \vfill
  % Título proyecto y escudo informática ----------------
  \colorbox{cpardoBox}{%
    \begin{minipage}{.8\textwidth}
      \vspace{.5cm}\Large
      \begin{center}
      \textbf{TFG del Grado en Ingeniería Informática}\vspace{.6cm}\\
      \textbf{\LARGE\@title{}}
      \end{center}
      \vspace{.2cm}
    \end{minipage}

  }%
  \hfill\begin{minipage}{.20\textwidth}
    \includegraphics[width=\textwidth]{escudoInfor}
  \end{minipage}
  \vfill
  % Datos de alumno, curso y tutores ------------------
  \begin{center}%
  {%
    \noindent\LARGE
    Presentado por \@author{}\\ 
    en Universidad de Burgos --- \@date{}\\
    Tutor: \@tutor{}\\
  }%
  \end{center}%
  \null
  \cleardoublepage
  }
\makeatother

\newcommand{\nombre}{Javier Martín Castro} %%% cambio de comando
\newcommand{\dni}{24392279C} %%% cambio de comando
\newcommand{\nomtutor}{Pedro Renedo Fernández} %%% cambio de comando
\newcommand{\nomcotutor}{} %%% cambio de comando

% Datos de portada
\title{Frontend aplicación móvil ayuda Esclerósis Múltiple}
\author{\nombre}
\tutor{\nomtutor}

\date{\today}

\begin{document}

\maketitle


\newpage\null\thispagestyle{empty}\newpage


%%%%%%%%%%%%%%%%%%%%%%%%%%%%%%%%%%%%%%%%%%%%%%%%%%%%%%%%%%%%%%%%%%%%%%%%%%%%%%%%%%%%%%%%
\thispagestyle{empty}


\noindent\includegraphics[width=\textwidth]{cabecera}\vspace{1cm}

\noindent D. \tutor, profesor del departamento de nombre departamento, área de nombre área.

\noindent Expone:

\noindent Que el alumno D. \nombre, con DNI  \dni, ha realizado el Trabajo final de Grado en Ingeniería Informática titulado título de TFG. 

\noindent Y que dicho trabajo ha sido realizado por el alumno bajo la dirección del que suscribe, en virtud de lo cual se autoriza su presentación y defensa.

\begin{center} %\large
En Burgos, {\large \today}
\end{center}

\vfill\vfill\vfill

% Author and supervisor
\begin{minipage}{0.45\textwidth}
\begin{flushleft} %\large
Vº. Bº. del Tutor:\\[2cm]
D. {\nomtutor}\\
\end{flushleft}
\end{minipage}
\hfill
\begin{minipage}{0.45\textwidth}
\begin{flushleft} %\large
Vº. Bº. del co-tutor:\\[2cm]
D. {\nomcotutor}\\
\end{flushleft}
\end{minipage}
\hfill

\vfill

% para casos con solo un tutor comentar lo anterior
% y descomentar lo siguiente
%Vº. Bº. del Tutor:\\[2cm]
%D. nombre tutor


\newpage\null\thispagestyle{empty}\newpage




\frontmatter

% Abstract en castellano
\renewcommand*\abstractname{Resumen}
\begin{abstract}
Dar solución a personas con \textbf{Esclerósis Múltiple} creando la interfaz visual.
\end{abstract}

\renewcommand*\abstractname{Descriptores}
\begin{abstract}
Palabras separadas por comas que identifiquen el contenido del proyecto Ej: aplicación móvil, android, ayuda, Esclerósis múltiple \ldots
\end{abstract}

\clearpage

% Abstract en inglés
\renewcommand*\abstractname{Abstract}
\begin{abstract}
A \textbf{brief} presentation of the topic addressed in the project.
\end{abstract}

\renewcommand*\abstractname{Keywords}
\begin{abstract}
keywords separated by commas.
\end{abstract}

\clearpage

% Indices
\tableofcontents

\clearpage

\listoffigures

\clearpage

\listoftables
\clearpage

\mainmatter
\capitulo{1}{Introducción}

El proyecto nace para dar solución y mejorar la calidad de vida de personas que tengan Esclerosis Múltiple (EM). Para ello se va a desarrollar el Frontend para dispositivo móvil Android. 

\capitulo{2}{Objetivos del proyecto}

% Este apartado explica de forma precisa y concisa cuales son los objetivos que se persiguen con la realización del proyecto. Se puede distinguir entre los objetivos marcados por los requisitos del software a construir y los objetivos de carácter técnico que plantea a la hora de llevar a la práctica el proyecto.

Crear una aplicación móvil para ayudar a personas con Esclerosis Múltiple, donde se realizarán las siguientes acciones:
\begin{itemize}
	\item Crear pantalla para acceder a la aplicación el usuario.
	\item Crear pantalla para realizar el registro de un nuevo usuario.
	\item Crear una pantalla del detalle del usuario.
	\item Crear pantalla home.
\end{itemize}

En el transcurso de la realización del proyecto se pretende aportar:
\begin{itemize}
	\item Dar soporte a personas con EM.
	\item Facilitar el día a día de los usuarios que usen la aplicación.
\end{itemize}



\capitulo{3}{Conceptos teóricos}

En aquellos proyectos que necesiten para su comprensión y desarrollo de unos conceptos teóricos de una determinada materia o de un determinado dominio de conocimiento, debe existir un apartado que sintetice dichos conceptos.

Algunos conceptos teóricos de \LaTeX{} \footnote{Créditos a los proyectos de Álvaro López Cantero: Configurador de Presupuestos y Roberto Izquierdo Amo: PLQuiz}.

\section{Secciones}

Las secciones se incluyen con el comando section.

\subsection{Subsecciones}

Además de secciones tenemos subsecciones.

\subsubsection{Subsubsecciones}

Y subsecciones. 


\section{Referencias}

Las referencias se incluyen en el texto usando cite~\cite{wiki:latex}. Para citar webs, artículos o libros~\cite{koza92}, si se desean citar más de uno en el mismo lugar~\cite{bortolot2005, koza92}.


\section{Imágenes}

Se pueden incluir imágenes con los comandos standard de \LaTeX, pero esta plantilla dispone de comandos propios como por ejemplo el siguiente:

\imagen{escudoInfor}{Autómata para una expresión vacía}{.5}



\section{Listas de items}

Existen tres posibilidades:

\begin{itemize}
	\item primer item.
	\item segundo item.
\end{itemize}

\begin{enumerate}
	\item primer item.
	\item segundo item.
\end{enumerate}

\begin{description}
	\item[Primer item] más información sobre el primer item.
	\item[Segundo item] más información sobre el segundo item.
\end{description}
	
\begin{itemize}
\item 
\end{itemize}

\section{Tablas}

Igualmente se pueden usar los comandos específicos de \LaTeX o bien usar alguno de los comandos de la plantilla.

\tablaSmall{Herramientas y tecnologías utilizadas en cada parte del proyecto}{l c c c c}{herramientasportipodeuso}
{ \multicolumn{1}{l}{Herramientas} & App AngularJS & API REST & BD & Memoria \\}{ 
HTML5 & X & & &\\
CSS3 & X & & &\\
BOOTSTRAP & X & & &\\
JavaScript & X & & &\\
AngularJS & X & & &\\
Bower & X & & &\\
PHP & & X & &\\
Karma + Jasmine & X & & &\\
Slim framework & & X & &\\
Idiorm & & X & &\\
Composer & & X & &\\
JSON & X & X & &\\
PhpStorm & X & X & &\\
MySQL & & & X &\\
PhpMyAdmin & & & X &\\
Git + BitBucket & X & X & X & X\\
Mik\TeX{} & & & & X\\
\TeX{}Maker & & & & X\\
Astah & & & & X\\
Balsamiq Mockups & X & & &\\
VersionOne & X & X & X & X\\
} 

\capitulo{4}{Técnicas y herramientas}
\section{Metodologías}
\subsection{Scrum}
Scrum~\cite{wiki:scrum} es un marco que se emplea para la gestión de proyectos de metodología ágil y que facilita a los equipos la tarea de gestionar y estructurar el trabajo. Para ello, a través de iteraciones (sprints) se realiza una serie de tareas en un periodo de tiempo y se van incorporando al software.
\subsection{Gitflow}
Gitflow~\cite{gitflow} se define como un flujo de trabajo para la creación de ramas en Git y llevar un control de versiones. En este flujo existen dos ramas principales, master y develop, y otras ramas que pueden ayudarnos en nuestro desarrollo como pueden ser: feature, release o  hotfix.
\section{Patrón de diseño}
\subsection{Model-View-ViewModel MVVM}
Es un patrón MVVM~\cite{mvvm}, también conocido como Model View ViewModel que se centra en separar la interfaz del usuario (View) de la parte lógica (Model). La interacción entre la parte lógica y la interfaz del usuario a través del recurso ViewModel.

Algunas de las ventajas al usar este patrón son:
\begin{itemize}
\item Fácil desarrollo ya que al poder separar la vista de la lógica varios equipos pueden trabajar simultáneamente en varios componentes.
\item Fácil testeo ya que no es necesario utilizar la vista para crear tests para el model o el viewmodel.
\item Fácil mantenimiento ya que al realizar la separación de los componentes se crea un código simple y limpio.
\end{itemize}

Vamos a describir a continuación cada uno de los componentes que forman el MVVM:
\begin{itemize}
	\item Model: es el componente donde se encapsulan los datos de nuestra aplicación. En ella se pueden encontrar validación y lógica de negocio.
	\item View: nos muestra el diseño y la apariencia de nuestra aplicación. En ella se verán los datos pero sin contener nada de la lógica de negocio.
	\item ViewModel: es el componente que enlaza los datos o cambios de estado que puede tener nuestra aplicación.
\end{itemize}

\imagen{mvvm-pattern}{Arquitectura limpia}{.9}

\section{Repositorio}
Entre las herramientas consideradas GitHub~\cite{github}, Bitbucket~\cite{bitbucket}, GitLab~\cite{wiki:gitlab} se decide utilizar GitHub, porque nos permite alojar proyectos gratuitamente. Además podemos crear documentación a través de wikis, crear tareas, sprints.
Al crear el respositorio en GitHub podemos usar el gestor de proyecto que tiene a través de sus paneles y saber el estado en el que está cada issue (tarea).
Otra de las ventajas que encontramos con GitHub es que está integrado en multiples servicios de integración continua.
\section{Control de versiones}
Al haber escogido como repositorio GitHub para alojar nuestro proyecto, utilizamos como control de versiones Git~\cite{wiki:git} que es un sistema distribuido que nos permite realizar nuestro trabajo sin conexión, es ligero y rápido.
El crear ramas y hacer los merges es muy rápido y pocas veces se obtienen conflictos.
El historial es muy detallado.
\section{Gestión del proyecto}
Como gestor de proyectos utilizamos el que viene integrado en GitHub. Con esta herramienta podemos gestionar nuestro proyecto y ver en que estado está cada tarea, incluso podemos darle prioridad.
\section{Entorno de desarrollo integrado (IDE)}
\subsection{Flutter}
Flutter~\cite{flutter} es un SDK (Kit de Desarrollo de Software) desarrollado por Google que nos permite crear aplicaciones para varios dispositivos. Hemos elegido esta herramienta porque nos permite crear aplicaciones multiplataforma.
Entre alguna de las ventajas destacamos:
\begin{itemize}
\item \textbf{Compila en nativo} tanto para Android como iOS.
\item Mediante el uso de widgets, se pueden \textbf{crear interfaces gráficas flexibles}.
\item \textbf{Rapidez} en el desarrollo.
\end{itemize}
\subsection{Dart}
Dart~\cite{dart} es un \textbf{lenguaje de programación moderno}, desarrollado por Google, donde se mezcla la programación orientada a objetos POO con los lenguajes de programación basados en scripts.
\subsection{FlutterFlow}
FlutterFlow~\cite{flutterflow} es una herramienta que nos permite crear aplicaciones insertando los componentes que tiene la herramienta, arrastrando y soltando. Es posible añadir lógica durante la creación de las páginas.
\subsection{LaTeX}
Entre las distintas herramientas de \LaTeX{} que existen hemos decidido utilizar TeXmaker~\cite{wiki:texmaker, texmaker} que es un editor de \LaTeX{} gratuito y multiplataforma. Además, incluye soporte unicode, corrección ortográfica, autocompletado y un visor de PDF integrado.
\subsection{Documentación}
Para la documentación del trabajo de fin de grado hemos empleado \LaTeX{}~\cite{wiki:latex} que es un sistema que permite componer textos con una alta calidad tipográfica. En el caso de la documentación del repositorio hemos empleado Markdown~\cite{wiki:markdown} que es un lenguaje de marcado ligero.
\section{Integración y entregas continuas (CI/CD)}
\subsection{Codemagic}
Codemagic~\cite{codemagic} es producto basado en la nube que permite realizar integración y entregas continuas. Te permite crear workflow versátiles y activar la compilación,  análisis de código y pruebas de manera automática.
\section{Herramientas contrucción del proyecto}
\subsection{Librerías}




\capitulo{5}{Aspectos relevantes del desarrollo del proyecto}

Este apartado pretende recoger los aspectos más relevantes del desarrollo del proyecto, pasando por la creación del proyecto, que decisiones se tomaron y la resolución de 
problemas encontrados y como se solucionaron.

\section{Inicio del proyecto}
El proyecto surgió a partir de la necesidad de ayudar a personas diagnosticadas con EM (Esclerósis Múltiple) y poder hacer su día a día más fácil a través de consejos, atención, ayuda.

Para poder hacer posible esta solución, se crea una aplicación inicialmente para el móvil, donde se incluirán secciones que pueden ayudar a los usuarios que hagan uso de ellas. Tras la asignación del proyecto se inicia un estudio de las posibilidades tecnológicas disponibles para resolver esta necesidad.

Desde el proyecto realizaremos llamadas a una API del cliente que nos proveerá de los datos necesarios para crear una aplicación amigable para los usuarios.

\section{Metodologías}
Desde el comienzo se realiza una arquitectura limpia para reutilizar los componentes y que sean independientes sin que haya ningún impacto si se realizan cambios en alguno de ellos.
Para la gestión del proyecto, se utiliza el panel que nos proporciona GitHub, usando la metodología ágil Scrum. El equipo al ser un proyecto creado para la realización del trabajo de fin de carrera, está formado por una persona. Por otro lado, al tener que hacer uso de una API externa, realizada por un compañero de la universidad, se mantienen reuniones para recabar requisitos técnicos para desarrollar la aplicación.

Para el desarrollo de la metodología ágil se siguieron los siguientes puntos:
\begin{itemize}
\item crear iteraciones incrementales (sprints) y revisiones
\item una vez terminado el sprint se añade a la rama main que será nuestro producto final.
\item los sprints por regla general los hemos hecho de duración semanal.,
\item para cada sprint se le asocian las tareas más importantes.
\end{itemize}

\section{Formación}
Para abordar la realización y creación del proyecto se necesitaban adquirir una serie de conocimientos que inicialmente eran inexistentes. Entre los conocimientos adquiridos se destacan:
\begin{itemize}
\item Formación sobre Flutter~\cite{flutter,formacionFlutter}.
\item Formación sobre Dart~\cite{dart}.
\end{itemize}

Además, se profundiza en obtener información a cerca de que es la esclerosis múltiple y que tipos existen.

\section{Desarrollo}
Durante el desarrollo de la aplicación nos aparecieron diversos retos:
\begin{itemize}
\item Conexión entre los diferentes componentes.
\item Internacionalización de la aplicación.
\item Contactar con el cliente ante desconexión de la dirección de la API facilitada inicialmente.
\item Contactar con el cliente para resolución de dudas sobre la API.
\end{itemize}

\section{Publicación}
Una vez hubo una versión lista para desplegar en producción, se creó una cuenta en Google Play y se desplegó.

Para próximas versiones, se tratará de crear un despliegue para otros dispositivos, como es iOS.




\capitulo{6}{Trabajos relacionados}

Este apartado sería parecido a un estado del arte de una tesis o tesina. En un trabajo final grado no parece obligada su presencia, aunque se puede dejar a juicio del tutor el incluir un pequeño resumen comentado de los trabajos y proyectos ya realizados en el campo del proyecto en curso. 

\capitulo{7}{Conclusiones y Líneas de trabajo futuras}

Todo proyecto debe incluir las conclusiones que se derivan de su desarrollo. Éstas pueden ser de diferente índole, dependiendo de la tipología del proyecto, pero normalmente van a estar presentes un conjunto de conclusiones relacionadas con los resultados del proyecto y un conjunto de conclusiones técnicas. 
Además, resulta muy útil realizar un informe crítico indicando cómo se puede mejorar el proyecto, o cómo se puede continuar trabajando en la línea del proyecto realizado. 



\bibliographystyle{plain}
\bibliography{bibliografia}

\end{document}
